\begin{document}

\lhead[ \leftmark ]{Informatik}
\rhead[Informatik]{Thomas Graf, \today, Winterthur}

\thispagestyle{empty}

\section{Klassentag}\label{sec:tag}

Liebe Schülerinnen und Schüler der Klasse 1c, liebe Eltern. Am Dienstag, 10. September 2024 wird der Klassentag der Klasse 1c stattfinden. Das Verschiebedatum bei schlechtem Wetter wäre der Freitag, 13. September 2024. 

Ziel des Klassentags ist die gemeinsame Planung des Klassenlagers (siehe \cref{sec:lager}). Die Schülerinnen und Schüler werden Menüpläne und Einkaufslisten ausarbeiten und Abendunterhaltungen konzipieren. 
Zudem bietet der Klassentag eine gute Gelegenheit, sich untereinander besser kennenzulernen und auszutauschen.

Am Mittag werden wir zusammen in den Wald oben beim Goldenberg gehen, ein Feuer anzünden und ein individuell mitgebrachtes Mittagessen einnehmen (Würste, Sandwiches, anderes Grillgut, $\ldots$).

\begin{description}
  \item[Grillstelle 1. Wahl:] \url{https://maps.app.goo.gl/3p5Uh1vmMrhA3anR9}
  \item[Grillstelle 2. Wahl:] \url{https://maps.app.goo.gl/dsiK2Exe1Hv661Eb6}
\end{description}

Abgesehen von dem Mittagessen wird der Klassentag im Schulhaus stattfinden.



\section{Klassenlager}\label{sec:lager}
Das Klassenlager wird am Montag, 30. September 2024 beginnen (Abfahrt ab Winterthur HB: 07:09 Uhr) und am Freitag, 04. Oktober 2024 enden (Ankunft im Winterthur HB: 16:51 Uhr). Unsere Unterkunft wird das \textit{Ferienhaus Madulain} im Engadin sein: \url{https://www.ferienhaus-madulain.ch/unterkunft}.

Das Klassenlager werde ich zusammen mit Stefanie Bäurle (Deutschlehrerin der Klasse 1c) leiten.

Weitere Informationen zum Klassenlager werden Sie am Elternabend (Donnerstag, 05. September 2024) erfahren und auch schriftlich erhalten.


\vspace{10cm}
Freundliche Grüsse\\
Thomas Graf, Klassenlehrer der Klasse 1c

\end{document}